% Curs 12 - Pagina 2 (3-I-1995).
% Fragment pentru % Curs 12 - Pagina 2 (3-I-1995).
% Fragment pentru % Curs 12 - Pagina 2 (3-I-1995).
% Fragment pentru % Curs 12 - Pagina 2 (3-I-1995).
% Fragment pentru \input{courses/pages/curs12-pagina2-ro}.
% Transcrierea respecta succesiunea si notatiile manuscrisului.
% Fara preambul si fara salt de pagina fortat.

\subsection*{Curs 12 - Pagina 2}

\noindent\underline{Ex.}\quad
$(\mathbb{Z},\cdot,1)\xrightarrow{\ f\ }(\mathbb{Z}_n,\cdot,\hat 1)$.

\noindent $f(e)=\hat e$.

\bigskip
\noindent
$(M_2(\mathbb{Z}),\cdot,I_2)\longrightarrow(\mathbb{Z},\cdot,1)$,
\qquad $A\longmapsto\det A$.
% In manuscris apar initial R la M_2(R) si la codomeniu, dar sub fiecare
% este scris Z (aparenta corectura de la R la Z). Am redat Z; de verificat
% pe original inainte de editia definitiva.

\bigskip
\noindent\underline{Obs.}\quad
Exist\u{a} \c{s}i func\c{t}ii care au (1), dar nu au (2)!
% In manuscris: "(\exists) si fct. care au (1) dar nu au (2)!".

\bigskip
\noindent $\ker f=R_f=\{(x,y)\in M\times M\mid f(x)=f(y)\}$.

\noindent $R_f$ e congruen\c{t}\u{a} pe $M$.

\bigskip
\noindent\underline{Dem.}\quad $R_f$ e rela\c{t}ie de echivalen\c{t}\u{a}$\,\sim$.

\medskip
\noindent $x\,R_f\,y\ \Longrightarrow$
\[
\begin{cases}
 zx\,R_f\,zy,\\
 xz\,R_f\,yz,
\end{cases}
\qquad (\forall)\ z\in M.
\]

\noindent
$f(xz)=f(x)f(z)=f(y)f(z)=f(yz)\ \Longrightarrow\ xz\,R_f\,yz$.

\medskip
\noindent $\Longrightarrow R_f$ e compatibil\u{a} cu ``$\cdot$''.
% Manuscrisul verifica explicit compatibilitatea la dreapta;
% compatibilitatea la stanga este enuntata, nu demonstrata pe rand separat.

\bigskip
\noindent\fbox{T.}\quad
\underline{De factorizare a morfismelor de monoizi}

\noindent (T. fundamental\u{a} a izomorfismelor de monoizi pentru morfisme de monoizi.)

\bigskip
\noindent Fie $f:M\longrightarrow M'$, $f$ morfism de monoizi.

\noindent $R\subseteq M\times M$ congruen\c{t}\u{a} pe $M$ a.\^{\i}.\ $R\subseteq R_f$.

\bigskip
\noindent\textbf{(1)}
\[
\begin{tikzcd}[column sep=large,row sep=large]
M \arrow[r,"f"] \arrow[d,"\varphi_R"'] & M' \\
M/R \arrow[ur,dashed,"f^{*}"'] &
\end{tikzcd}
\]
\noindent $\Longrightarrow\ (\exists!)\ f^{*}:M/R\longrightarrow M'$ a.\^{\i}.\
$f^{*}\circ\varphi_R=f$.
% Proiectia catre M/R este indicata in manuscris prin f_R (sau phi_R);
% aici \varphi_R este citirea provizorie a simbolului greu de distins.

\bigskip
\noindent\textbf{(2)}\quad $f^{*}$ inj. $\Longleftrightarrow R=R_f$.

\medskip
\noindent\textbf{(3)}\quad $f^{*}$ surj. $\Longleftrightarrow f$ surj.

\medskip
\noindent $f^{*}$ bij. $\Longleftrightarrow$\,\footnote{\^{I}n manuscris, o s\u{a}geat\u{a} curb\u{a} leag\u{a} aceast\u{a} echivalen\c{t}\u{a} de condi\c{t}iile (2) \c{s}i (3), f\u{a}r\u{a} a repeta termenul din dreapta.}

\bigskip
\noindent\fbox{T.}\quad (fundamental\u{a} de izomorfism).\quad
$f:M\longrightarrow M'$ morfism de monoizi
\[
\Longrightarrow\quad \operatorname{Im} f\cong\frac{M}{R_f}.
\]
\noindent\hfill $f(M)\leq M'$.
% Textul "f(M) <= M'" este scris sub "Im f" in original.
% Pagina se incheie aici.
.
% Transcrierea respecta succesiunea si notatiile manuscrisului.
% Fara preambul si fara salt de pagina fortat.

\subsection*{Curs 12 - Pagina 2}

\noindent\underline{Ex.}\quad
$(\mathbb{Z},\cdot,1)\xrightarrow{\ f\ }(\mathbb{Z}_n,\cdot,\hat 1)$.

\noindent $f(e)=\hat e$.

\bigskip
\noindent
$(M_2(\mathbb{Z}),\cdot,I_2)\longrightarrow(\mathbb{Z},\cdot,1)$,
\qquad $A\longmapsto\det A$.
% In manuscris apar initial R la M_2(R) si la codomeniu, dar sub fiecare
% este scris Z (aparenta corectura de la R la Z). Am redat Z; de verificat
% pe original inainte de editia definitiva.

\bigskip
\noindent\underline{Obs.}\quad
Exist\u{a} \c{s}i func\c{t}ii care au (1), dar nu au (2)!
% In manuscris: "(\exists) si fct. care au (1) dar nu au (2)!".

\bigskip
\noindent $\ker f=R_f=\{(x,y)\in M\times M\mid f(x)=f(y)\}$.

\noindent $R_f$ e congruen\c{t}\u{a} pe $M$.

\bigskip
\noindent\underline{Dem.}\quad $R_f$ e rela\c{t}ie de echivalen\c{t}\u{a}$\,\sim$.

\medskip
\noindent $x\,R_f\,y\ \Longrightarrow$
\[
\begin{cases}
 zx\,R_f\,zy,\\
 xz\,R_f\,yz,
\end{cases}
\qquad (\forall)\ z\in M.
\]

\noindent
$f(xz)=f(x)f(z)=f(y)f(z)=f(yz)\ \Longrightarrow\ xz\,R_f\,yz$.

\medskip
\noindent $\Longrightarrow R_f$ e compatibil\u{a} cu ``$\cdot$''.
% Manuscrisul verifica explicit compatibilitatea la dreapta;
% compatibilitatea la stanga este enuntata, nu demonstrata pe rand separat.

\bigskip
\noindent\fbox{T.}\quad
\underline{De factorizare a morfismelor de monoizi}

\noindent (T. fundamental\u{a} a izomorfismelor de monoizi pentru morfisme de monoizi.)

\bigskip
\noindent Fie $f:M\longrightarrow M'$, $f$ morfism de monoizi.

\noindent $R\subseteq M\times M$ congruen\c{t}\u{a} pe $M$ a.\^{\i}.\ $R\subseteq R_f$.

\bigskip
\noindent\textbf{(1)}
\[
\begin{tikzcd}[column sep=large,row sep=large]
M \arrow[r,"f"] \arrow[d,"\varphi_R"'] & M' \\
M/R \arrow[ur,dashed,"f^{*}"'] &
\end{tikzcd}
\]
\noindent $\Longrightarrow\ (\exists!)\ f^{*}:M/R\longrightarrow M'$ a.\^{\i}.\
$f^{*}\circ\varphi_R=f$.
% Proiectia catre M/R este indicata in manuscris prin f_R (sau phi_R);
% aici \varphi_R este citirea provizorie a simbolului greu de distins.

\bigskip
\noindent\textbf{(2)}\quad $f^{*}$ inj. $\Longleftrightarrow R=R_f$.

\medskip
\noindent\textbf{(3)}\quad $f^{*}$ surj. $\Longleftrightarrow f$ surj.

\medskip
\noindent $f^{*}$ bij. $\Longleftrightarrow$\,\footnote{\^{I}n manuscris, o s\u{a}geat\u{a} curb\u{a} leag\u{a} aceast\u{a} echivalen\c{t}\u{a} de condi\c{t}iile (2) \c{s}i (3), f\u{a}r\u{a} a repeta termenul din dreapta.}

\bigskip
\noindent\fbox{T.}\quad (fundamental\u{a} de izomorfism).\quad
$f:M\longrightarrow M'$ morfism de monoizi
\[
\Longrightarrow\quad \operatorname{Im} f\cong\frac{M}{R_f}.
\]
\noindent\hfill $f(M)\leq M'$.
% Textul "f(M) <= M'" este scris sub "Im f" in original.
% Pagina se incheie aici.
.
% Transcrierea respecta succesiunea si notatiile manuscrisului.
% Fara preambul si fara salt de pagina fortat.

\subsection*{Curs 12 - Pagina 2}

\noindent\underline{Ex.}\quad
$(\mathbb{Z},\cdot,1)\xrightarrow{\ f\ }(\mathbb{Z}_n,\cdot,\hat 1)$.

\noindent $f(e)=\hat e$.

\bigskip
\noindent
$(M_2(\mathbb{Z}),\cdot,I_2)\longrightarrow(\mathbb{Z},\cdot,1)$,
\qquad $A\longmapsto\det A$.
% In manuscris apar initial R la M_2(R) si la codomeniu, dar sub fiecare
% este scris Z (aparenta corectura de la R la Z). Am redat Z; de verificat
% pe original inainte de editia definitiva.

\bigskip
\noindent\underline{Obs.}\quad
Exist\u{a} \c{s}i func\c{t}ii care au (1), dar nu au (2)!
% In manuscris: "(\exists) si fct. care au (1) dar nu au (2)!".

\bigskip
\noindent $\ker f=R_f=\{(x,y)\in M\times M\mid f(x)=f(y)\}$.

\noindent $R_f$ e congruen\c{t}\u{a} pe $M$.

\bigskip
\noindent\underline{Dem.}\quad $R_f$ e rela\c{t}ie de echivalen\c{t}\u{a}$\,\sim$.

\medskip
\noindent $x\,R_f\,y\ \Longrightarrow$
\[
\begin{cases}
 zx\,R_f\,zy,\\
 xz\,R_f\,yz,
\end{cases}
\qquad (\forall)\ z\in M.
\]

\noindent
$f(xz)=f(x)f(z)=f(y)f(z)=f(yz)\ \Longrightarrow\ xz\,R_f\,yz$.

\medskip
\noindent $\Longrightarrow R_f$ e compatibil\u{a} cu ``$\cdot$''.
% Manuscrisul verifica explicit compatibilitatea la dreapta;
% compatibilitatea la stanga este enuntata, nu demonstrata pe rand separat.

\bigskip
\noindent\fbox{T.}\quad
\underline{De factorizare a morfismelor de monoizi}

\noindent (T. fundamental\u{a} a izomorfismelor de monoizi pentru morfisme de monoizi.)

\bigskip
\noindent Fie $f:M\longrightarrow M'$, $f$ morfism de monoizi.

\noindent $R\subseteq M\times M$ congruen\c{t}\u{a} pe $M$ a.\^{\i}.\ $R\subseteq R_f$.

\bigskip
\noindent\textbf{(1)}
\[
\begin{tikzcd}[column sep=large,row sep=large]
M \arrow[r,"f"] \arrow[d,"\varphi_R"'] & M' \\
M/R \arrow[ur,dashed,"f^{*}"'] &
\end{tikzcd}
\]
\noindent $\Longrightarrow\ (\exists!)\ f^{*}:M/R\longrightarrow M'$ a.\^{\i}.\
$f^{*}\circ\varphi_R=f$.
% Proiectia catre M/R este indicata in manuscris prin f_R (sau phi_R);
% aici \varphi_R este citirea provizorie a simbolului greu de distins.

\bigskip
\noindent\textbf{(2)}\quad $f^{*}$ inj. $\Longleftrightarrow R=R_f$.

\medskip
\noindent\textbf{(3)}\quad $f^{*}$ surj. $\Longleftrightarrow f$ surj.

\medskip
\noindent $f^{*}$ bij. $\Longleftrightarrow$\,\footnote{\^{I}n manuscris, o s\u{a}geat\u{a} curb\u{a} leag\u{a} aceast\u{a} echivalen\c{t}\u{a} de condi\c{t}iile (2) \c{s}i (3), f\u{a}r\u{a} a repeta termenul din dreapta.}

\bigskip
\noindent\fbox{T.}\quad (fundamental\u{a} de izomorfism).\quad
$f:M\longrightarrow M'$ morfism de monoizi
\[
\Longrightarrow\quad \operatorname{Im} f\cong\frac{M}{R_f}.
\]
\noindent\hfill $f(M)\leq M'$.
% Textul "f(M) <= M'" este scris sub "Im f" in original.
% Pagina se incheie aici.
.
% Transcrierea respecta succesiunea si notatiile manuscrisului.
% Fara preambul si fara salt de pagina fortat.

\subsection*{Curs 12 - Pagina 2}

\noindent\underline{Ex.}\quad
$(\mathbb{Z},\cdot,1)\xrightarrow{\ f\ }(\mathbb{Z}_n,\cdot,\hat 1)$.

\noindent $f(e)=\hat e$.

\bigskip
\noindent
$(M_2(\mathbb{Z}),\cdot,I_2)\longrightarrow(\mathbb{Z},\cdot,1)$,
\qquad $A\longmapsto\det A$.
% In manuscris apar initial R la M_2(R) si la codomeniu, dar sub fiecare
% este scris Z (aparenta corectura de la R la Z). Am redat Z; de verificat
% pe original inainte de editia definitiva.

\bigskip
\noindent\underline{Obs.}\quad
Exist\u{a} \c{s}i func\c{t}ii care au (1), dar nu au (2)!
% In manuscris: "(\exists) si fct. care au (1) dar nu au (2)!".

\bigskip
\noindent $\ker f=R_f=\{(x,y)\in M\times M\mid f(x)=f(y)\}$.

\noindent $R_f$ e congruen\c{t}\u{a} pe $M$.

\bigskip
\noindent\underline{Dem.}\quad $R_f$ e rela\c{t}ie de echivalen\c{t}\u{a}$\,\sim$.

\medskip
\noindent $x\,R_f\,y\ \Longrightarrow$
\[
\begin{cases}
 zx\,R_f\,zy,\\
 xz\,R_f\,yz,
\end{cases}
\qquad (\forall)\ z\in M.
\]

\noindent
$f(xz)=f(x)f(z)=f(y)f(z)=f(yz)\ \Longrightarrow\ xz\,R_f\,yz$.

\medskip
\noindent $\Longrightarrow R_f$ e compatibil\u{a} cu ``$\cdot$''.
% Manuscrisul verifica explicit compatibilitatea la dreapta;
% compatibilitatea la stanga este enuntata, nu demonstrata pe rand separat.

\bigskip
\noindent\fbox{T.}\quad
\underline{De factorizare a morfismelor de monoizi}

\noindent (T. fundamental\u{a} a izomorfismelor de monoizi pentru morfisme de monoizi.)

\bigskip
\noindent Fie $f:M\longrightarrow M'$, $f$ morfism de monoizi.

\noindent $R\subseteq M\times M$ congruen\c{t}\u{a} pe $M$ a.\^{\i}.\ $R\subseteq R_f$.

\bigskip
\noindent\textbf{(1)}
\[
\begin{tikzcd}[column sep=large,row sep=large]
M \arrow[r,"f"] \arrow[d,"\varphi_R"'] & M' \\
M/R \arrow[ur,dashed,"f^{*}"'] &
\end{tikzcd}
\]
\noindent $\Longrightarrow\ (\exists!)\ f^{*}:M/R\longrightarrow M'$ a.\^{\i}.\
$f^{*}\circ\varphi_R=f$.
% Proiectia catre M/R este indicata in manuscris prin f_R (sau phi_R);
% aici \varphi_R este citirea provizorie a simbolului greu de distins.

\bigskip
\noindent\textbf{(2)}\quad $f^{*}$ inj. $\Longleftrightarrow R=R_f$.

\medskip
\noindent\textbf{(3)}\quad $f^{*}$ surj. $\Longleftrightarrow f$ surj.

\medskip
\noindent $f^{*}$ bij. $\Longleftrightarrow$\,\footnote{\^{I}n manuscris, o s\u{a}geat\u{a} curb\u{a} leag\u{a} aceast\u{a} echivalen\c{t}\u{a} de condi\c{t}iile (2) \c{s}i (3), f\u{a}r\u{a} a repeta termenul din dreapta.}

\bigskip
\noindent\fbox{T.}\quad (fundamental\u{a} de izomorfism).\quad
$f:M\longrightarrow M'$ morfism de monoizi
\[
\Longrightarrow\quad \operatorname{Im} f\cong\frac{M}{R_f}.
\]
\noindent\hfill $f(M)\leq M'$.
% Textul "f(M) <= M'" este scris sub "Im f" in original.
% Pagina se incheie aici.
